\section*{Chapitre \ref{ch:g3:stages-of-life} --- Les étapes de la vie}

\begin{solution}{exo:g3:stages-of-life:1}
Bébé, enfant, \omterm{def:g3:stages-of-life:stages}{adolescent}, \omterm{def:g3:stages-of-life:stages}{adulte}, \omterm{def:g3:stages-of-life:stages}{vieillesse}.
\end{solution}

\begin{solution}{exo:g3:stages-of-life:2}
(a) Pendant l'enfance, à partir de six ans environ. (b) À l'étape
\omterm{def:g3:stages-of-life:stages}{adulte}. (c) À l'étape du bébé --- la première année est la plus
rapide de toutes.
\end{solution}

\begin{solution}{exo:g3:stages-of-life:3}
La poussée du bébé et celle de l'\omterm{def:g3:stages-of-life:stages}{adolescent}, avec entre les deux la
croissance calme et régulière de l'enfance.
\end{solution}

\begin{solution}{exo:g3:stages-of-life:4}
La ligne devient plate~: l'âge continue d'avancer en bas, mais la
taille ne monte plus.
\end{solution}

\begin{solution}{exo:g3:stages-of-life:5}
Entre $6$ et $8$ ans~: d'environ $\qty{115}{cm}$ à environ
$\qty{127}{cm}$ --- à peu près $\qty{12}{cm}$ sur les deux années.
Pendant la seule première année~: d'environ $\qty{50}{cm}$ à environ
$\qty{75}{cm}$ --- à peu près $\qty{25}{cm}$, deux fois plus que ces
deux années d'enfance réunies.
\end{solution}

\begin{solution}{exo:g3:stages-of-life:6}
Parce qu'une comparaison n'est honnête que si les mesures sont
faites de la même façon. Des chaussures un jour et des chaussettes
le lendemain, ou un livre penché n'importe comment, créeraient une
«~croissance~» (ou la cacheraient) qui n'a jamais eu lieu.
\end{solution}

\begin{solution}{exo:g3:stages-of-life:7}
Gland (la \omterm{def:g2:from-seed-to-plant:seed}{graine}), \omterm{prop:g2:from-seed-to-plant:cycle}{jeune pousse}, jeune chêne, chêne \omterm{def:g3:stages-of-life:stages}{adulte} qui
fleurit et porte des glands, vieux chêne --- et, avec le temps, la
mort.
\end{solution}

\begin{solution}{exo:g3:stages-of-life:8}
Réponse libre. Un enfant de neuf ans est à l'étape enfant~; la
poussée de l'adolescence est encore devant lui.
\end{solution}

\begin{solution}{exo:g3:stages-of-life:9}
La mouche (quelques semaines), le chat (un an environ), l'humain
(près de vingt ans), le chêne (plusieurs dizaines d'années). Chez
tous les quatre, l'ordre des étapes est le même et ne peut être ni
sauté ni parcouru à l'envers~: jeune, \omterm{def:g3:stages-of-life:stages}{adulte}, \omterm{def:g3:stages-of-life:stages}{vieillesse}.
\end{solution}

\begin{solution}{exo:g3:stages-of-life:10}
Non. À neuf ans, il est dans la longue période calme de l'enfance
--- environ $\qty{6}{cm}$ par \emph{an}, c'est-à-dire un demi-centimètre
par mois~: bien trop peu pour se voir entre deux mesures espacées
d'un mois. La courbe monte sûrement, mais lentement~; il lui suffit
d'attendre quelques mois entre deux points.
\end{solution}
